\section{Boundary corrections and integrable metrics}
\label{g:bt}

Integration along a fibration can lose regularity along its discriminant.
The correction needed to extend the resulting metric also determines which
sections can be returned to the original variety.  We keep these two roles
together.  Three outputs will be useful: a metric integrable after
multiplication by a boundary section; an unweighted integrable metric; and
a nef anti-log-canonical class for a signed boundary.  Their slope
consequences are proved under the hypotheses actually available.

A metric weight is normalized by $|e|^2=\exp(-\psi)$ in a holomorphic
frame.  Integrability at exponent one therefore means integrability of
$\exp(-\psi)$.  For a rational line bundle we clear denominators and
divide the resulting weight by the denominator.  A signed rational
divisor $B$ on a smooth variety defines a \emph{sub-klt pair} if every
coefficient of its crepant boundary on a log resolution is less than one.
Negative coefficients are allowed; a sub-klt divisor need not be an
effective klt boundary.

\subsection{Approximation without a strict curvature hypothesis}
\label{g:bt:approximation}

The approximation below uses a fixed ample auxiliary line to produce
sections.  The metric being approximated is only semipositive.  We give
both the ideal-theoretic construction and the Gaussian construction,
since the latter directly preserves the weighted integral.

\begin{lemma}[Weighted sub-klt approximation]
\label{g:bt:weighted-approximation}
Let $Y$ be smooth projective, of dimension $d>0$, let $J$ be a rational
line bundle with a semipositive metric of weights $\psi$, and let $P\geq0$
be a rational divisor.  Suppose that, in local divisor frames,
\[
                 |s_P|^2e^{-\psi}\in L^1_{\mathrm{loc}}.
\]
There is a fixed ample Cartier divisor $A$ such that, for every sufficiently
large divisible integer $m$, there is an effective rational divisor
$\Delta_m\sim_{\mathbb Q}J+A/m$ for which $(Y,\Delta_m-P)$ is sub-klt.
When $P=0$, the pair $(Y,\Delta_m)$ is klt.
\end{lemma}

\begin{proof}
Choose a very ample $H$ and an ample $A$ such that $A-K_Y-iH$ is ample for
$1\leq i\leq d$.  Nadel vanishing and Castelnuovo--Mumford regularity give
global generation of
\[
       \mathcal O_Y(mJ+A)\otimes\mathcal I(m\psi).
\]
Indeed its twist by $-iH$ is an adjoint sheaf with coefficient metric of
strictly positive curvature, supplied by $A-K_Y-iH$, and multiplier ideal
$\mathcal I(m\psi)$.  The ideal is coherent and nonzero; local point
extension where $\psi$ is finite gives a nonzero integrable germ
\cite{Nadel,OhsawaTakegoshi,raDemailly}.

Let $u:Y'\to Y$ resolve this ideal and $P$, and write
$\mathcal I(m\psi)\mathcal O_{Y'}=\mathcal O_{Y'}(-F_m)$.
For a prime $E$ on $Y'$ put
\[
 b_E=\operatorname{ord}_E F_m,\qquad
 p_E=\operatorname{ord}_E u^*P,\qquad
 a_E=\operatorname{ord}_E K_{Y'/Y}.
\]
We claim that
\begin{equation}\label{g:bt:strict-approximation}
                         b_E/m-p_E<1+a_E.
\end{equation}
The local point-extension estimate for $m\psi$ supplies holomorphic
functions $v_x$ with $|v_x(x)|^2=e^{m\psi(x)}$ and uniformly bounded
$L^2(e^{-m\psi})$ norm on fixed nested coordinate balls.  The local upper
bound for $\psi$ bounds their ordinary norms, hence their suprema on a
smaller ball.  Each $v_x$ belongs to $\mathcal I(m\psi)$.  At a general
point of $E=(z_1=0)$, Cauchy estimates after division by $z_1^{b_E}$ give
\[
               e^{m\psi\circ u}\leq C_m|z_1|^{2b_E}.
\]
The change-of-variables formula bounds the pulled-back integrable density
below by a positive constant times $|z_1|^{2(p_E+a_E-b_E/m)}$.
Its transverse integral is finite only if
\eqref{g:bt:strict-approximation} holds.

Global generation identifies the fixed divisor of the pulled-back
system with $F_m$.  Choose a general member of the free residual system
and divide the corresponding divisor downstairs by $m$.  The crepant
boundary of $\Delta_m-P$ has coefficient $b_E/m-p_E-a_E<1$ on each fixed
component, and coefficient $1/m<1$ on new moving components.  Bertini,
applied also to the strata of the resolved divisor, makes their union
simple normal crossings.  Decreasing coefficients from the positive
part preserves positive log discrepancies, proving sub-klt on every
higher model.

There is a second way to choose the member.  The global Bergman lower
estimate for the same sufficiently ample $A$ gives local expressions
$s_{i,m}$ of an orthonormal basis with
\[
                   \sum_i|s_{i,m}|^2\geq C_m e^{m\psi}.
\]
This globalization by a fixed ample twist is the point-extension
construction in \cite[Section~9, proof of Proposition~9.1]{Demailly1993}.
One obtains the estimate by extending from a point, cutting off, and
solving $\bar\partial$ globally with a logarithmic pole of order $d$ at
that point.  The fixed positivity of $A-K_Y$ pays the uniformly bounded
negative curvature of the cutoff pole; integrability forces the
correction to vanish at the point.  Singular weights are handled by
decreasing regularization and weak limits in the $L^2$ estimate
\cite[Corollary~13.9, Remark~13.10, and Theorem~14.2]{raDemailly}.
For independent standard complex Gaussians $c_i$ and $m>1$,
\[
 \mathbb E\left|\sum_i c_i s_{i,m}\right|^{-2/m}
   =C'_m\left(\sum_i|s_{i,m}|^2\right)^{-1/m},
\]
where $C'_m<\infty$ because $\mathbb E|c_1|^{-2/m}<\infty$.
Multiplication by $|s_P|^2$ and Fubini show that almost every such member
has the weighted integrability required for
\eqref{g:bt:strict-approximation}.
\end{proof}

The approximation retains the negative part $-P$.  Its immediate use is
a cotangent estimate.  We first isolate the common slope argument.

\begin{lemma}[From foliation determinants to tensor bounds]
\label{g:bt:slope-principle}
Let $Y$ be smooth projective and $P\geq0$ rational.  Suppose that
$-K_Y+P$ is pseudoeffective and that $K_{\mathcal F}-K_Y+P$ is
pseudoeffective for every algebraically integrable saturated foliation
$\mathcal F\subset T_Y$ of positive rank and corank.  Then
$hP-c_1(M)$ is pseudoeffective for every line subsheaf
$M\subseteq(\Omega_Y^1)^{\otimes h}$, $h\geq0$.
\end{lemma}

\begin{proof}
The assertion on a point is immediate. For positive dimension,
fix a nonzero movable class $\alpha$ and put $t=\alpha\cdot P\geq0$.
The total degree of $T_Y$ is at least $-t$.  If its minimum
Harder--Narasimhan slope were less than $-t$, its positive-slope part
$\mathcal F$ would be nonzero and proper: the last negative graded piece
already has degree less than $-t$, whereas the whole sheaf does not.
The quotient by the positive part has only nonpositive slopes and degree
less than $-t$.  In particular
\[
 \mu_{\alpha,\min}(\mathcal F)>0,\qquad
 2\mu_{\alpha,\min}(\mathcal F)>
       \mu_{\alpha,\max}(T_Y/\mathcal F).
\]
The bracket map vanishes by the slope inequality, and the
Campana--P\u aun algebraicity theorem makes $\mathcal F$ algebraically
integrable \cite[Proposition~2.2]{Ou2023}.  But
$c_1(T_Y/\mathcal F)=K_{\mathcal F}-K_Y$ has degree less than $-t$,
contradicting the hypothesis.  Thus
$\mu_{\alpha,\min}(T_Y)\geq-t$.
Duality and the tensor-product rule for movable slopes give
$c_1(M)\cdot\alpha\leq ht$ \cite{GKP2016,CampanaPaun2019}.
Movable-cone duality proves the assertion \cite{BDPP2013}.
For $h=0$, the inclusion in $\mathcal O_Y$ proves it directly.
\end{proof}

\begin{proposition}[Finite-volume cotangent estimate]
\label{g:bt:weighted-slope}\label{g:slope-recalled}\label{rb:cotangent}
Let $Y$ be smooth projective and $P\geq0$ rational.  If $J=-K_Y+P$
has a semipositive metric with $|s_P|^2e^{-\psi}$ locally integrable, then
$hP-c_1(M)$ is pseudoeffective for every line subsheaf
$M\subseteq(\Omega_Y^1)^{\otimes h}$.
If $\dim Y>0$ and $\alpha\ne0$ is movable, then
$\mu_{\alpha,\min}(T_Y)\geq-P\cdot\alpha$. In particular,
$P\cdot\alpha=0$ gives $\mu_{\alpha,\min}(T_Y)\geq0$ and
$c_1(M)\cdot\alpha\leq0$.
\end{proposition}

\begin{proof}
The assertion on a point is immediate. Suppose $\dim Y>0$.
Lemma~\ref{g:bt:weighted-approximation} gives effective
$\Delta_m\sim_{\mathbb Q}-K_Y+P+A/m$ with
$(Y,\Delta_m-P)$ sub-klt.  Resolve a first-integral map for an algebraic
foliation $\mathcal F$ together with that pair, by $u:Y'\to Y$.
Write its crepant boundary as $B_m^+-B_m^-$, with disjoint effective
parts.  Then
\[
 K_{Y'}+B_m^+\sim_{\mathbb Q}u^*A/m+B_m^-,
 \qquad 0\leq u_*B_m^-\leq P.
\]
The positive boundary is klt.  A general fibre of the resolved
first-integral map meets the isomorphism locus of $u$; hence $u^*A$
restricts to a big divisor there.  A divisible positive multiple of
$K_{Y'}+B_m^+$ consequently has a nonzero section on that fibre.
Druel's relative pseudoeffectivity theorem, in the form recalled by Ou,
applies to the effective boundary and its log canonical general fibre
\cite[Proposition~4.1]{Druel2017},
\cite[Proposition~3.1]{Ou2023}. It yields pseudoeffectivity of
$K_{\mathcal F'}+B_m^+$.  Push down, add $P-u_*B_m^-\geq0$, and let $m$
tend to infinity.  This gives $K_{\mathcal F}-K_Y+P$ pseudoeffective.
The metric already gives $-K_Y+P$ pseudoeffective, so
Lemma~\ref{g:bt:slope-principle} applies.  This is the signed finite-volume
adaptation of Ou's tangent-slope argument; its relative positivity input
is due to Druel.
\end{proof}

\subsection{An unweighted integral and an annihilating movable class}
\label{g:bt:unweighted}

When the integral contains no boundary factor, orbifold cotangent
semipositivity gives a different estimate.  Its hypothesis concerns
$K_Y+J$, which need not be supported on an allowed-pole divisor.

\begin{proposition}[Unweighted movable-slope estimate]
\label{g:bt:unweighted-slope}
Let $Y$ be smooth projective of dimension $d>0$, and let $J$ have a
semipositive metric with $e^{-\psi}\in L^1_{\mathrm{loc}}$.
Suppose $K_Y+J$ is pseudoeffective.  For every movable class $\alpha$
and every line subsheaf $M\subseteq(\Omega_Y^1)^{\otimes h}$, $h>0$,
\begin{equation}\label{g:bt:orbifold-bound}
 c_1(M)\cdot\alpha\leq
           h d^{h-1}(K_Y+J)\cdot\alpha.
\end{equation}
Thus $\alpha\cdot(K_Y+J)=0$ gives $c_1(M)\cdot\alpha\leq0$.
For $h=0$ the latter inequality holds without the annihilation premise.
\end{proposition}

\begin{proof}
Use Lemma~\ref{g:bt:weighted-approximation} with $P=0$ to obtain klt
$\Delta_m\sim_{\mathbb Q}J+A/m$.  On a log resolution $u:Y'\to Y$
truncate the negative coefficients of the crepant boundary.  This gives
an effective simple normal crossing boundary $\Delta'_m$, with
coefficients less than one, and
\[
 K_{Y'}+\Delta'_m=u^*(K_Y+\Delta_m)+E_m,\qquad
 E_m\geq0\quad\hbox{$u$-exceptional}.
\]
The left side is pseudoeffective.  The numerical pullback $\alpha'$,
defined by $\alpha'\cdot D=\alpha\cdot u_*D$, is movable and annihilates
exceptional divisors.

On an adapted finite cover $v:\widehat Y\to Y'$, the orbifold cotangent
sheaf $\mathcal E$ contains $v^*\Omega^1_{Y'}$ and satisfies
$c_1(\mathcal E)=v^*(K_{Y'}+\Delta'_m)$.  In a boundary chart its
generators are pullbacks of $dx/x^\delta$, where $\delta$ is the boundary
coefficient and the cover clears its denominator.  Campana--P\u aun
tensor-quotient semipositivity gives nonnegative $v^*\alpha'$-degree to
torsion-free quotients of $\mathcal E^{\otimes h}$
\cite[Theorem~1.3]{CampanaPaun2019}.
The line $v^*u^*M$ injects through ordinary cotangent pullback.
Saturating increases its determinant by an effective divisor.  The
quotient has nonnegative degree; if its rank is zero, the same inequality
follows directly from the effective saturation divisor.  Since
$c_1(\mathcal E^{\otimes h})=h d^{h-1}c_1(\mathcal E)$, cancellation of
the finite covering degree gives
\[
 c_1(M)\cdot\alpha\leq
    h d^{h-1}(K_Y+J+A/m)\cdot\alpha.
\]
Let $m$ increase to obtain \eqref{g:bt:orbifold-bound}.
\end{proof}

For $J=-K_Y+\Theta$ with $\Theta\geq0$, cone duality gives
$h d^{h-1}\Theta-c_1(M)$ pseudoeffective.
For $J=-K_Y+kN+R$, the proposition instead applies to a movable class
annihilating $N$ and $R$.  Both situations occur below.

\subsection{Weighted transfer with arbitrary exceptional poles}
\label{g:bt:exceptional}

Let $X$ be smooth projective, let $L=-K_X$ carry a smooth semipositive
metric, and let a torus $T$ act with the natural linearization.  Average
the metric over a compact real form.  For a $T$-invariant rational map to
a smooth projective $Y$, let
\[
                 X\xleftarrow{\pi}Z\xrightarrow{g}Y
\]
be a smooth equivariant resolution, with geometrically integral generic
fibre $F$ over $K=\mathbb C(Y)$.  We distinguish the following strong
rank condition: there is a nonzero invariant section $s$ of
$a\pi^*L|_F$, $a>0$, and
\begin{equation}\label{g:bt:all-exceptional-rank}
 \dim_K H^0\bigl(F,(b\pi^*L+D)|_F\bigr)^T\leq1
\end{equation}
for every integer $b\geq0$ and every effective invariant
$\pi$-exceptional integral divisor $D$.
This condition persists on higher equivariant resolutions: push down
sections and test regularity at codimension one.  Exceptional divisors
for an additional resolution remain exceptional on the generic fibre.

Flattening by modification, followed by finite normalization, provides
a normal projective model $p:W\to X$, $f:W\to Y$ with no prime divisor
mapping into codimension at least two in $Y$ \cite{RaynaudGruson1971}.
Here and below the base function field is unchanged.
Regard $s$ as a rational section on $W$.  For every base prime $Q$, set
\[
 e_E=\operatorname{ord}_E f^*Q,\qquad
 n_Q=\min_{E\mapsto Q}\frac{\operatorname{ord}_E s}{e_E},
 \qquad N=\sum_Q n_Q Q.
\]
The sums have finite support.  Let $B$ be the reduced sum of $Q$ for
which every component above $Q$ is exceptional over $X$.  Away from $B$
define $n'_Q$ by the same minimum restricted to nonexceptional components;
on $B$ put $n'_Q=n_Q$.  Set $R=\sum_Q(n'_Q-n_Q)Q\geq0$.
Taking a power of $s$, make $N,R$ integral.  Then
\begin{equation}\label{g:bt:normalization}
 D_s=\operatorname{div}_W(s)-f^*N\geq0,\qquad
       f^*\mathcal O_Y(N)\longrightarrow a p^*L
       \quad\hbox{is regular}.
\end{equation}
Normality and the divisorial image property justify regularity.

\begin{lemma}[Extension of the normalized maximum]
\label{g:bt:bounded-maximum}
In the normalization \eqref{g:bt:normalization}, suppose the negative
logarithms of the adjoint moments are psh for all sufficiently large
exponents on a smooth connected-fibre open set.
Then the maximum norm gives a semipositive metric with locally bounded
weights on $N$.
\end{lemma}

\begin{proof}
In a local frame of $N$, write $\sigma$ for the regular normalized
section.  On the smooth open set, the moments have limit
\[
 -m^{-1}\log\int_{Z_y}|\sigma|^{2m}
                \frac{\pi^*d\mu}{g^*d\lambda_Y}
                     \longrightarrow-\log\max_{Z_y}|\sigma|^2.
\]
The fibre measure is finite and has dense positive locus.
After normalization to total mass one, its power means increase to
the maximum.  Proper smoothness makes that maximum continuous.
Integration over a fixed patch on which both factors are nonzero
gives local upper bounds for the rescaled psh weights.  The limit is
therefore psh.  No positivity of the zeroth moment is needed.

Regularity of $\sigma$ on the proper model gives a local lower bound
for the limiting weight up to the boundary.
Above each base prime a component attaining the vertical minimum has
$\sigma\ne0$ at its general point.  In transverse coordinates the map
is $y_1=x_1^{e_E}$, with the other base coordinates free.  A neighbourhood
there meets every nearby fibre and bounds its maximum away from zero.
This gives the upper bound needed for psh extension at general points
of every omitted divisor.  Psh extension across the remaining
codimension-two set then applies; local upper boundedness there follows
from transverse discs with boundary off the removed set.  The lower
bound persists.  The extended weights are locally bounded and obey the
frame transformation rule for $N$.
\end{proof}

\begin{proposition}[Weighted exceptional correction]
\label{g:bt:weighted-transfer}
Under \eqref{g:bt:all-exceptional-rank}, $N$ has a semipositive metric
with locally bounded weights.  For some integer $c\geq0$,
$-K_Y+cB$ has a semipositive metric satisfying
$|s_B|^{2c}e^{-\psi}\in L^1_{\mathrm{loc}}$.
A nonzero section of any positive divisible multiple of
$j(N+R)+lB$, with $j>0$ and $l\geq0$, gives a nonzero invariant section
of a positive multiple of $L$ on $X$.
\end{proposition}

\begin{proof}
If $Y$ is a point, all the base divisors vanish and their metrics are
positive constants. The section $s$ is then regular on $Z$ and pushes
down to the required section on $X$. Assume henceforth that $\dim Y>0$.

For every $m\geq0$, \eqref{g:bt:all-exceptional-rank} applied to
$D=K_{Z/X}$ gives invariant rank one for
$g_*(K_{Z/Y}+(ma+1)\pi^*L)$.  Its generator is the product of $s^m$,
the invariant Jacobian section, and a local inverse canonical frame of
$Y$.  In a local frame of $N$, let $\sigma$ be the normalized section
from \eqref{g:bt:normalization}.  The generator has squared norm
\[
 I_m(y)=\int_{Z_y}|\sigma|^{2m}\,
                         \frac{\pi^*d\mu}{g^*d\lambda_Y},
\]
where $d\mu$ is the smooth volume defined by the metric of $L$.
Lemma~\ref{g:moments} applies: Berndtsson positivity is used on the
smooth base-change locus and the invariant line is a holomorphic
orthogonal summand.  Lemma~\ref{g:bt:bounded-maximum} then gives the
bounded metric of $N$.

For the second metric use $\rho=I_0$.  Its negative logarithm is psh on
the smooth locus, and $\rho\,d\lambda_Y=g_*(\pi^*d\mu)$ has finite mass.
At a general point of a prime $E\mapsto Q$, choose transverse coordinates
with $y_1=x_1^{e_E}$ and the other base coordinates free.
If $a_E=\operatorname{ord}_E K_{Z/X}$, integration on a small patch gives
\[
                 \rho(y)\geq C|y_1|^{2((1+a_E)/e_E-1)}.
\]
Outside $B$ choose a nonexceptional component, so $a_E=0$.
On the finitely many components of $B$ choose $c$ large enough that
\[
                     \psi=-\log\rho+c\log|s_B|^2
\]
is locally bounded above at their general points.  The same bound is
already available at every other omitted divisor.  Psh extension across
these divisors, then across the remaining set of codimension at least
two, gives the metric on $-K_Y+cB$.  Its weighted integral is exactly
the original finite pushforward mass.

Finally, represent a section of a positive divisible multiple of
$j(N+R)+lB$ by a rational function $v$ on $Y$. After incorporating
that multiple in $j,l$, its divisor satisfies
\[
             \operatorname{div}(v)+j(N+R)+lB\geq0.
\]
Return the rational section $s^jf^*v$. At a nonexceptional component
$E$ above $Q\not\subset B$, its order is at least
$j(h_E-e_En'_Q)\geq0$, where $h_E=\operatorname{ord}_Es$.
There are no nonexceptional components above $B$, and
equidimensionality excludes vertical primes over higher-codimension
sets. Horizontal orders are nonnegative by the generic regularity of
$s$. Thus every pole is exceptional over $X$ and disappears on
pushdown. All factors are invariant.
\end{proof}

\subsection{Unweighted toroidal transfer}
\label{g:bt:toroidal}

The preceding correction uses only a divisorially equidimensional model.
A toroidal model gives a finer rational correction and an unweighted
integral.  The rank hypothesis in the next proposition includes zero.

\begin{proposition}[Integrable toroidal correction]
\label{g:bt:toroidal-transfer}
Use the smooth metric and action above.  Let $p:W\to X$ be birational
with $W$ normal, and let $f:W\to Y$ be projective, equidimensional and
toroidal for strict toroidal embeddings, with $Y$ smooth.  Include the
exceptional locus of $p$ in the source boundary.  Suppose that an
invariant rational section $s$ of $aL$, $a>0$, is regular on the generic
fibre and that on a smooth equivariant resolved diagram
\begin{equation}\label{g:bt:zeroth-rank}
 \dim_K H^0\bigl(F,(K_{Z/X}+ma\pi^*L)|_F\bigr)^T=1
                           \qquad(m\geq0).
\end{equation}
Normalize $s$ by vertical minima as in \eqref{g:bt:normalization}.
For a boundary prime $Q$ put
\[
 A_E=1+\operatorname{ord}_E K_{W/X},\qquad
 c_Q=\min_{E\mapsto Q}\frac{A_E}{e_E},\qquad
 \Theta=\sum_Q\max\{c_Q-1,0\}Q.
\]
Then $N$ has a semipositive metric with locally bounded weights, and
$-K_Y+\Theta$ has a semipositive metric with locally integrable squared
frame norms.  The pullback $f^*\Theta$ is exceptional over $X$.
\end{proposition}

\begin{proof}
For a point base, $N=\Theta=0$ and both metrics are positive constants.
We may therefore assume $\dim Y>0$.

Lemma~\ref{g:moments} supplies the psh moment weights, and
Lemma~\ref{g:bt:bounded-maximum} gives the metric of $N$.
The zeroth moment is the pushed-forward smooth volume density
$\rho$.  It can be computed on a toroidal resolution of $W$, even when
that model is not equivariant: change of variables identifies its
density with the one on an equivariant resolution.
Put $\theta_Q=\max\{c_Q-1,0\}$ and use
\[
                      \psi=-\log\rho+
                                   \sum_Q\theta_Q\log|s_Q|^2.
\]
At a minimizing component the transverse calculation in the preceding
proof gives $\rho\geq C|s_Q|^{2(c_Q-1)}$.  This is the local upper
bound needed to extend $\psi$ at general divisor points.
Off the boundary the morphism is smooth and gives a positive lower
bound.  Extend over the remaining codimension-two set by psh removability.

The correction $\Theta$ is the divisor $A_0$ in
Lemma~\ref{g:volume-subklt}. To compare the two definitions, it remains
only to check primes outside the base boundary. A vertical prime above
such a prime cannot belong to the source boundary of a toroidal
morphism. Since that boundary contains the exceptional locus, the
prime is nonexceptional over $X$; the morphism is smooth at its general
point. Its Jacobian order is zero and its pullback multiplicity is one,
so the corresponding coefficient of $A_0$ is zero.

The rank-free discrepancy lemma now gives both the exceptional support
of $f^*\Theta$ and, on every smooth model $Z'\to W$ with maps
$\pi':Z'\to X$ and $g':Z'\to Y$,
\[
       \operatorname{coeff}_F(K_{Z'/X}-(g')^*\Theta)>-1
                         \qquad\text{for every prime }F\subset Z'.
\]
On a simultaneous normal-crossing resolution, these are precisely the
exponents of the corrected source measure in the projection formula
\[
 e^{-\psi}\,d\lambda_Y
   =g'_*\left(\prod_Q|s_Q\circ g'|^{-2\theta_Q}
                                  (\pi')^*d\mu\right).
\]
They make that measure locally integrable. Proper pushforward proves
the required unweighted integrability on $Y$, including at boundary
intersections. Curvature here came from the invariant zeroth moment;
the discrepancy argument used no adjoint-rank assumption.
\end{proof}

No horizontal strict-curvature patch was assumed here.
The output is the integrability used in
Proposition~\ref{g:bt:unweighted-slope}.  The separate toroidal-volume
theorem includes the additional patch needed by its effectivity
applications.

\subsection{The large-moment construction for eventual rank one}
\label{g:bt:eventual}

We now construct the donor metric from a positive moment. The hypothesis
will initially specify adjoint rank only at sufficiently large degrees.
Let $F_X=\mathbb C(X)$, let $K\subseteq F_X^T$ be finitely generated and
relatively algebraically closed in $F_X$, and let $s$ be a nonzero
invariant rational section of $dL$, $d>0$.
For a prime divisorial valuation $E$ over $X$, write
$a_E=\operatorname{ord}_E K_{X'/X}$ on a smooth model carrying $E$.
Assume
\begin{equation}\label{g:bt:horizontal-regular}
 \operatorname{ord}_E s\geq0
       \quad\hbox{whenever $\operatorname{ord}_E|_{K^*}=0$},
\end{equation}
and, for every sufficiently large integer $m$,
\begin{equation}\label{g:bt:eventual-rank}
 \left\{t\in\operatorname{Rat}(mdL)^T:
   \begin{array}{l}
    \operatorname{ord}_E t+a_E\geq0\\[-2pt]
    \text{for all $E$ trivial on $K$}
   \end{array}\right\}=K s^m.
\end{equation}
Here $\operatorname{Rat}(mdL)$ is the one-dimensional $F_X$-space of
rational sections, with zero included in the displayed set.
Both conditions concern all divisorial valuations.  On one smooth
resolution of the map to a model of $K$, they can be tested at
horizontal primes: the divisor inequalities pull to further resolutions
because the new relative canonical divisor is effective.

Together these two conditions imply rank one at every nonnegative
multiple of $d$, including zero. Indeed, if $t$ has degree $jd$, with
$j\geq0$ an integer, and satisfies
the displayed adjoint inequalities, choose $m\geq j$ in the eventual
range. Horizontal regularity of $s$ makes $ts^{m-j}$ satisfy the same
inequalities at degree $md$. Equation~\eqref{g:bt:eventual-rank} gives
$ts^{m-j}\in Ks^m$, hence $t\in Ks^j$; the reverse inclusion follows
from regularity of $s$ and $a_E\geq0$. The distinction of the following
construction is its metric: it uses one positive moment and retains
the corresponding multiple of the normalized section line.

Choose a model $p:W\to X$, $f:W\to Y$ with $Y$ smooth projective,
$W$ normal and $\mathbb Q$-Gorenstein klt, and every vertical prime
mapping onto a base prime.  Such a model is obtained by toroidalization
and subdivision.  Refine target cones by source-cone images; on the
induced source subdivision each ray maps to a ray or to zero.
The reduced boundary of the initial smooth toroidal source defines a
positive slicing function on its cone complex, linear on each cone and
positive away from the origin. Its pullback remains
positive on the induced subdivision, including on horizontal rays.
Abramovich--Rojas \cite[Corollary~0.11]{AbramovichRojas1998} then gives
a coherent simplicial subdivision without new rays. Since the target
cones are smooth, every new source cone still maps to a target face;
equidimensionality is preserved. The resulting toroidal source has
quotient singularities, hence is klt. These are birational modifications,
with no field extension \cite{ADK2013,AK2000,KKMSD}.

For each base prime $Q$, put
\[
 h_E=\operatorname{ord}_E s,\qquad e_E=\operatorname{ord}_E f^*Q,\qquad
 n_Q=\min_{E\mapsto Q}h_E/e_E,\qquad N=\sum_Qn_Q Q,
\]
and $D=\operatorname{div}_W(s)-f^*N\geq0$.  Define
\[
 r_Q=\min_{\substack{E\mapsto Q\\D_E=0}}a_E/e_E,\qquad
                       R=\sum_Qr_Q Q\geq0.
\]
The minimum is over a nonempty finite set.  Only finitely many
coefficients of $N,R$ are nonzero.

\begin{proposition}[An integrable large moment]
\label{g:bt:large-moment}
There is an integer $k$ in the range of \eqref{g:bt:eventual-rank} such
that
\[
            \Xi=K_{W/X}+kD-f^*R\geq0.
\]
The rational line $N$ has a semipositive metric with locally bounded
weights, and $J=-K_Y+kN+R$ has a semipositive metric with
$\mathcal I(\psi)=\mathcal O_Y$.
A nonzero section of a positive divisible multiple of $N+R/k$ returns
to a nonzero invariant section of a positive multiple of $L$ on $X$.
\end{proposition}

\begin{proof}
At a component with $D_E=0$, the definition of $r_Q$ gives
$a_E-e_Er_Q\geq0$.  At the finitely many other relevant components,
increase $k$ to absorb $e_Er_Q$ by $kD_E$.  At horizontal components
there is no subtraction.  This proves $\Xi\geq0$, and over every $Q$
some component has $\Xi_E=0$.

On a smooth equivariant model over $Y$ define
\[
 I_m\,d\lambda_Y=g_*\bigl(|\pi^*s|^{2m}\pi^*d\mu\bigr).
\]
For all sufficiently large $m$, \eqref{g:bt:eventual-rank} identifies
the invariant adjoint line generated by the Jacobian times $s^m$.
Smooth direct-image positivity and orthogonal compact averaging imply
that $-\log I_m$ is psh on a common smooth open.
The densities are smooth there, so the base-change locus for an
individual $m$ can be removed by continuity.  High moments give
\[
 I_m^{1/m}\longrightarrow\beta(y)=\max_{Z_y}|\pi^*s|^2.
\]
For this limit one may divide by $I_0$ as a positive density on the
smooth open; no curvature assertion about $I_0$ is used.
The weights of the required lines are
\begin{equation}\label{g:bt:large-weights}
 \begin{split}
 \varphi_N&=-\log\beta+\sum_Q n_Q\log|s_Q|^2,\\
 \psi&=-\log I_k+\sum_Q(kn_Q+r_Q)\log|s_Q|^2.
 \end{split}
\end{equation}
The first extends with locally bounded weights by regularity of the
normalized section and a minimizing component above every divisor,
by Lemma~\ref{g:bt:bounded-maximum}.

For the second, choose a component above $Q$ with $\Xi_E=0$.
In coordinates $y_1=x_1^{e_E}$ at its general point, the corrected
source density has exponent
\[
          a_E+kh_E-e_E(kn_Q+r_Q)=0.
\]
Its pushforward is bounded below by $C|y_1|^{-2(1-1/e_E)}$.
Thus $\psi$ is bounded above at general points of every omitted
divisor and extends psh there, then across codimension two.
On a log resolution $v:\widehat W\to W$ the exponents in the corrected
source density are the coefficients of
\[
 K_{\widehat W/X}+kv^*D-v^*f^*R
                   =K_{\widehat W/W}+v^*\Xi.
\]
They are all greater than $-1$, since $W$ is klt and $\Xi$ is effective
and $\mathbb Q$-Cartier.  Normal crossing integration and properness
give $e^{-\psi}\in L^1_{\mathrm{loc}}$ on the entire base.

Finally $D-f^*R/k$ is nonnegative at each component nonexceptional
over $X$, since $a_E=0$ there and $\Xi\geq0$.  Pulling back a section
of $N+R/k$ and multiplying by a power of $s$ leaves only exceptional
poles, removed on pushdown.
\end{proof}

The resulting donor line is $J=-K_Y+kN+R$. Its $kN$ term records the
positive moment used to construct it. Although the hypotheses also
imply rank one at zero, this proof obtains integrability from $I_k$;
it does not produce the same metric on $-K_Y+R$.

\subsection{Signed anti-log-canonical classes}
\label{g:bt:signed}

The third slope input starts from a nef class rather than a weighted
integral.  It will be applied to the signed discriminant produced by the
toroidal density calculation.  Its negative part measures exactly the
possible loss in the cotangent estimate.

\begin{proposition}[Signed-boundary cotangent estimate]
\label{g:bt:signed-slope}
Let $Y$ be smooth projective and let $B$ be a rational divisor with simple
normal crossing support and coefficients less than one.  Write
$B=B^+-B^-$ in disjoint positive and negative parts, and suppose that
$-(K_Y+B)$ is nef.  Then
\[
                         hB^- -c_1(M)
                         \quad\hbox{is pseudoeffective}
\]
for every line subsheaf $M\subseteq(\Omega_Y^1)^{\otimes h}$, $h\geq0$.
If $\dim Y>0$ and $\alpha\ne0$ is a movable class with
$\alpha\cdot B^-=0$, this gives
$\mu_{\alpha,\min}(T_Y)\geq0$ and $c_1(M)\cdot\alpha\leq0$.
\end{proposition}

\begin{proof}
The class $-K_Y+B^-=-(K_Y+B)+B^+$ is pseudoeffective.
For an algebraically integrable foliation $\mathcal F$ of positive rank
and corank, resolve its rational first-integral map together with $B$ by
$u:Y'\to Y$, and write
\[
             K_{Y'}+C=u^*(K_Y+B),\qquad C=C^+-C^-.
\]
Every coefficient of $C$ is less than one: the positive part of the
original simple normal crossing boundary is klt, and decreasing its
coefficients can only increase discrepancies.  In particular
$(Y',C^+)$ is klt.  For an ample Cartier $H$ on $Y'$ and rational
$\delta>0$, choose a general effective rational divisor
\[
        A_\delta\sim_{\mathbb Q}
                    -u^*(K_Y+B)+\delta H
\]
such that $(Y',C^++A_\delta)$ is klt.  The class is ample, so a sufficiently
divisible multiple supplies such a member by Bertini, also on a general
fibre of the resolved map.  We have
\[
        K_{Y'}+C^++A_\delta\sim_{\mathbb Q} C^-+\delta H.
\]
It is big on that fibre and hence has a nonzero section after taking a
positive divisible multiple.  Apply Druel's relative pseudoeffectivity
theorem \cite[Proposition~4.1]{Druel2017}, in the form of
\cite[Proposition~3.1]{Ou2023}, to conclude that
$K_{\mathcal F'}+C^++A_\delta$ is pseudoeffective.
Its pushforward is numerically
\[
                 K_{\mathcal F}-K_Y+B^-+\delta u_*H.
\]
Letting $\delta$ decrease to zero proves
$K_{\mathcal F}-K_Y+B^-$ pseudoeffective.
Lemma~\ref{g:bt:slope-principle} now proves the assertion for every
movable class, and therefore in the pseudoeffective cone.
\end{proof}

The distinction between the three slope inputs is now explicit.
Proposition~\ref{g:bt:weighted-slope} retains the integral of a specified
boundary section.  Proposition~\ref{g:bt:unweighted-slope} uses an
unweighted integral and the class $K_Y+J$.
Proposition~\ref{g:bt:signed-slope} uses nefness of an anti-log-canonical
class with a signed boundary.  Each supplies the cotangent inequality
needed for the corresponding contraction; none asserts a global
strict-curvature lower bound.
